%==============================================================================
\part{The Language}
%==============================================================================

\section{The source model}
\label{sec:model}

\subsection{Sources differ in kind, not dialect}

The design turns on a fact about the data landscape that is easy to state and
consequential to accept.

\begin{observation}[Heterogeneity of kind]
\label{obs:kind}
Among the public resources a biological question routinely crosses, one finds
at least the following four kinds:
\begin{enumerate}[label=(\roman*),leftmargin=2.6em]
\item a SPARQL protocol endpoint accepting arbitrary graph patterns;
\item a resource exposing both a SPARQL endpoint and an independent REST
interface, with different coverage and different identifier conventions on the
two;
\item a REST interface returning flat delimited text, with a fixed finite set
of operations and no query language whatsoever;
\item an artefact --- an ontology file, a mapping table --- distributed for
local loading, whose query capability is whatever the loader provides.
\end{enumerate}
\end{observation}

Kind (iii) is the one that forecloses the obvious designs. A resource whose
entire query surface is a fixed family of request shapes such as
\texttt{/link/\{target\}/\{source\}} and \texttt{/get/\{id\}} does not admit a
compilation target for arbitrary patterns. It admits exactly the operations it
names. Any language that compiles to ``a query'' must therefore either exclude
such resources or define ``query'' so broadly that the word does no work. We
take the second horn and rename: what a source accepts is a \emph{request}, and
what shapes the request is a \emph{capability set}.

\begin{definition}[Capability vocabulary]
\label{def:capability}
Fix a finite vocabulary $\Feat$ of \emph{capabilities}. A capability names a
construct that a request may use. Throughout,
\[
\Feat = \{\textsf{pattern}, \textsf{path}, \textsf{bind}, \textsf{agg},
\textsf{filter}, \textsf{neg}, \textsf{order}, \textsf{lookup},
\textsf{link}, \textsf{batch}, \textsf{regex}\},
\]
where \textsf{pattern} is a conjunctive pattern over the source's vocabulary,
\textsf{path} a regular path expression, \textsf{bind} the inline supply of a
value set as input, \textsf{agg} an aggregate, \textsf{filter} a scalar
predicate evaluated at the source, \textsf{neg} negation, \textsf{order} a
sort, \textsf{lookup} retrieval of a record by primary identifier,
\textsf{link} retrieval of identifiers in a second namespace related to a given
identifier, \textsf{batch} the acceptance of more than one input identifier per
request, and \textsf{regex} a string match.
\end{definition}

The vocabulary is not canonical and nothing below depends on its particular
members; what the development requires is only that $\Feat$ be finite and that
capabilities occupy syntactically independent positions in each concrete
request language (\cref{asm:independence}).

\begin{definition}[Source]
\label{def:source}
A \emph{source} is a quadruple
$\Src = (\eta_\Src, \Capset(\Src), \Ext_\Src, \cost_\Src)$
where $\eta_\Src$ is an opaque endpoint handle, $\Capset(\Src) \subseteq \Feat$ is
the \emph{declared capability set}, $\Ext_\Src$ is the \emph{extraction
function} mapping a raw response to a result set, and
$\cost_\Src : \Nset \to \Rset_{>0}$ gives the cost of a request as a function
of its input cardinality.
\end{definition}

\begin{example}[Four sources]
\label{ex:sources}
A reaction knowledge base with a SPARQL endpoint declares
\[
\{\textsf{pattern},\textsf{path},\textsf{bind},\textsf{agg},\textsf{filter},
\textsf{neg},\textsf{order},\textsf{regex},\textsf{batch}\}.
\]
A flat-file pathway REST interface declares $\{\textsf{lookup},\textsf{link}\}$
and nothing else --- in particular not \textsf{filter}, so a plan wishing to
filter a pathway result must retrieve and filter locally, at a cost the plan
can see. A locally loaded ontology declares
$\{\textsf{pattern},\textsf{path},\textsf{filter}\}$ but not \textsf{batch},
since the loader admits one query at a time. A protein resource exposing both a
SPARQL endpoint and a REST interface is modelled as \emph{two} sources with
different capability sets, for the reason given in \cref{rem:not-union}.
\end{example}

\begin{remark}[A dual-interface resource is two sources]
\label{rem:not-union}
Modelling a resource with both SPARQL and REST access as a single source whose
capability set is the union of the two is unsound. The union asserts that some
request may use \textsf{path} and \textsf{link} together, and no request can:
the two capabilities belong to different interfaces, with different coverage
and often different identifier conventions. Capability sets are per-interface,
and a plan crossing the two interfaces of one resource crosses two sources and
pays two costs.
\end{remark}

\subsection{Results, and why the result model is the common one}

\begin{definition}[Result set]
\label{def:result}
A \emph{result set} over namespace $n$ is a finite set
\[
\Res \subseteq n \times (\mathcal{A} \rightharpoonup \mathcal{V})
\]
of pairs (identifier, partial attribute map), in which no identifier occurs
twice. We write $\idm(\Res) \subseteq n$ for the projection onto identifiers,
and $|\Res| = |\idm(\Res)|$.
\end{definition}

Making the identifier primary and the attributes secondary is the whole of the
portability argument, so it is worth establishing that the choice is available.

\begin{proposition}[The result model is common where the query model is not]
\label{prop:common-result}
Each of the four source kinds of \cref{obs:kind} admits an extraction function
into \cref{def:result} that is defined on every syntactically valid response.
\end{proposition}

\begin{proof}
For kind (i), and for the SPARQL interface of kind (ii), a response is a
solution sequence: a list of maps from projected variables to values. Designate
one projected variable as the identifier column and the remainder as
attributes; collapse duplicate identifiers by merging their attribute maps.
The construction is defined on every well-formed solution sequence with at
least one projected variable.

For the REST interfaces of kinds (ii) and (iii), a response is one of two
shapes. Under \textsf{lookup} it is a record keyed by the requested identifier,
yielding the singleton whose identifier is that key and whose attribute map is
the record's fields. Under \textsf{link} it is a delimited list of pairs,
yielding one element per line whose identifier is the second column and whose
attribute map records the first column as provenance. Both are defined on every
syntactically valid response.

For kind (iv), the loader's own result form is that of kind (i), and the
construction of the first paragraph applies. In each case the construction
consumes the entire response, so no information is discarded silently, though
information may be relocated into attributes.
\end{proof}

\begin{remark}[What \cref{prop:common-result} does not say]
\label{rem:extraction-lossy}
It does not say the extraction is lossless in the sense that matters. A
solution sequence with three projected variables and no functional dependency
on any one of them --- a genuine ternary relation --- is forced into
identifier-plus-attributes by choosing a key, and the choice loses the
multiplicity of the other two columns. A plan that needs an $n$-ary relation
not functionally determined by a key cannot express it in this model.
\Cref{sec:limits} records this as a real restriction rather than a formality,
and \cref{rem:nary-workaround} gives the standard, unsatisfying workaround.
\end{remark}

\begin{remark}[The workaround, and what it costs]
\label{rem:nary-workaround}
The standard response to the restriction of \cref{rem:extraction-lossy} is
reification: introduce a synthetic identifier for each tuple of the $n$-ary
relation and carry the original columns as attributes of it. The result is a
result set in the sense of \cref{def:result}, so the model accepts it, and
nothing is lost in the sense of information content.

What is lost is everything downstream. A synthetic identifier lies in no
namespace of any source, so it is outside the domain of every translation map
of \cref{def:map}; a reified step therefore cannot be the input of a
translation step, and the retention arithmetic of \cref{thm:retention} has
nothing to say about it. Nor can it be joined against a source on identity,
since no source will have minted the same synthetic key. In effect reification
buys expressiveness at the price of federation: the reified set is a terminal
value the plan may emit but may not carry across a boundary. We call the
workaround standard because it is what one does, and unsatisfying because a
language for federated retrieval that answers ``express it as something you
cannot federate'' has not answered.
\end{remark}

\begin{remark}[Disjoint namespaces are a modelling decision]
\label{rem:disjoint}
We assume namespaces are pairwise disjoint. In practice they are not: the same
lexical string may be a valid accession in two databases, and one database may
mint identifiers in another's namespace as aliases. The assumption is
discharged in the prototype by prefixing every identifier with its namespace
tag at extraction time, which makes disjointness true by construction at the
cost of making equality between an identifier and its own alias fail. That
failure is visible as a retention loss in \cref{sec:translation} rather than as
a silent join miss, which we regard as the better trade, but it is a trade.
\end{remark}

\subsection{Lowering, faithfulness, and honesty}

\begin{definition}[Abstract request]
\label{def:areq}
An \emph{abstract request} is a triple $\rho = (\phi, \Req(\rho), \beta)$ where
$\phi$ is a pattern over a declared vocabulary, $\Req(\rho) \subseteq \Feat$ is
the set of capabilities that $\phi$ uses, and $\beta$ is a possibly empty tuple
of result sets to be supplied as input. We write $\llbracket \rho \rrbracket_D$
for the \emph{denotation} of $\rho$ over a dataset $D$: the result set that
$\rho$ specifies, defined independently of any source.
\end{definition}

\begin{definition}[Lowering]
\label{def:lowering}
A \emph{lowering} for source $\Src$ is a partial function
$\Low_\Src : \rho \rightharpoonup \mathcal{W}_\Src$ into the source's concrete
request language. $\Low_\Src$ is \emph{faithful for} $f \in \Feat$ if for every
abstract request $\rho$ with $\Req(\rho) \subseteq \{f\} \cup \Capset(\Src)$ and
every dataset $D$ the source holds,
\[
\Ext_\Src\bigl(\eta_\Src(\Low_\Src(\rho))\bigr) \;=\; \llbracket \rho \rrbracket_D .
\]
\end{definition}

\begin{definition}[Honest declaration]
\label{def:honest}
$\Capset(\Src)$ is \emph{honest} if, for every $f \in \Feat$,
\[
f \in \Capset(\Src) \iff \Low_\Src \text{ is faithful for } f .
\]
\end{definition}

\begin{assumption}[Syntactic independence]
\label{asm:independence}
For each source $\Src$ there is an injection from $\Feat$ into a set of
pairwise disjoint syntactic positions of $\mathcal{W}_\Src$, such that the
concrete request realising a set of capabilities is the one filling exactly the
corresponding positions, and its denotation is the conjunction of the
denotations of the parts.
\end{assumption}

\Cref{asm:independence} holds for the concrete languages we target --- a
pattern, a path inside a pattern, a binding block, an aggregate wrapper, a
filter clause and a sort modifier are separate positions in SPARQL, and a REST
interface with a fixed operation family satisfies it degenerately because each
capability corresponds to a distinct operation. It fails for languages in which
two constructs interact semantically, and \cref{rem:independence-fails} records
one such interaction that is not hypothetical.

\begin{theorem}[Lowering is total and faithful exactly on contained requests]
\label{thm:lowering}
Let $\Src$ be a source whose declaration is honest and which satisfies
\cref{asm:independence}. Then for every abstract request $\rho$,
\[
\Low_\Src(\rho) \text{ is defined and faithful}
\quad\Longleftrightarrow\quad
\Req(\rho) \subseteq \Capset(\Src).
\]
\end{theorem}

\begin{proof}
($\Leftarrow$) Suppose $\Req(\rho) \subseteq \Capset(\Src)$. By
\cref{def:honest}, $\Low_\Src$ is faithful for each $f \in \Req(\rho)$
individually. By \cref{asm:independence} the capabilities in $\Req(\rho)$
occupy pairwise disjoint syntactic positions, so the concrete request filling
all of those positions is well formed and lies in $\mathcal{W}_\Src$; hence
$\Low_\Src(\rho)$ is defined. Again by \cref{asm:independence} its denotation
is the conjunction of the denotations of its parts, and by faithfulness for
each $f$ separately each part denotes what $\rho$'s corresponding component
specifies. The conjunction therefore equals $\llbracket \rho \rrbracket_D$, so
$\Low_\Src(\rho)$ is faithful.

($\Rightarrow$) Suppose $f \in \Req(\rho) \setminus \Capset(\Src)$ for some $f$.
By honesty, $\Low_\Src$ is not faithful for $f$: there exist an abstract
request $\rho'$ using $f$ and a dataset $D$ on which either $\Low_\Src(\rho')$
is undefined or the extracted response differs from
$\llbracket \rho' \rrbracket_D$. In the first case $\Low_\Src$ is undefined on
some request using $f$, and by \cref{asm:independence} the position that $f$
occupies is unfillable for $\Src$, hence $\Low_\Src(\rho)$ is undefined as
well. In the second case the position is fillable but its part denotes
something other than what $\rho$ specifies, and since by
\cref{asm:independence} the whole denotation is the conjunction over parts,
$\Low_\Src(\rho)$ is defined but unfaithful. Either way the left-hand side
fails.
\end{proof}

\begin{remark}[Honesty is an assumption, and the interesting one]
\label{rem:honesty-assumption}
\Cref{thm:lowering} is exactly as strong as \cref{def:honest}, and honesty is
asserted by the adapter author, not verified against the service. There are two
ways to violate it and they are not symmetric. \emph{Under}-declaration --- a
capability the source in fact supports faithfully, omitted from $\Capset(\Src)$ ---
is safe: it costs expressiveness, causes some plans to be refused that could
have been answered, and never produces a wrong answer. \emph{Over}-declaration
is unsound and invisible: the adapter claims \textsf{path}, the service accepts
path expressions but evaluates them under a different semantics or truncates
them, the lowering is defined, a response comes back, and it is not the
denotation. This is not hypothetical --- \cref{sec:interp} exhibits a case in
which a public service accepted a construct and returned a count differing from
the specified denotation by a factor of $198.5$. We therefore treat the
capability set as a claim to be minimised rather than a feature list to be
maximised, and record in \cref{sec:limits} that no mechanical check enforces
this, and that we know of no way to build one short of differential testing
against a reference implementation.
\end{remark}

\begin{proposition}[Under-declaration can be forced by the service, not the language]
\label{prop:under}
There exist a source $\Src$, a capability $f$, and a concrete language
$\mathcal{W}_\Src$ such that $\mathcal{W}_\Src$ supports $f$ in full
generality, $\Low_\Src$ emits well-formed requests using $f$, every such
request receives a well-formed response, and honesty nonetheless requires
$f \notin \Capset(\Src)$.
\end{proposition}

\begin{proof}
Take $f = \textsf{agg}$ and let $\Src$ be a SPARQL endpoint operating under a
server-side cap $M$ on the number of solutions it will materialise. Aggregation
is expressible in SPARQL in full generality, so $\mathcal{W}_\Src$ supports
$f$. Let $\Low_\Src$ emit a counting aggregate over a pattern $\phi$ whose
unaggregated solution sequence has cardinality $N > M$ on the dataset $D$ the
source holds. The endpoint materialises $M$ solutions, computes the aggregate
over those, and returns a well-formed response carrying the number $M$. But
$\llbracket \rho \rrbracket_D = N \neq M$. So $\Low_\Src$ is defined on $\rho$
and unfaithful for $f$, and by \cref{def:honest} honesty forces
$f \notin \Capset(\Src)$ even though the target language supports $f$ and the
request succeeded.
\end{proof}

\Cref{prop:under} is the reason \cref{def:honest} is phrased in terms of the
lowering and the service jointly rather than in terms of the target language's
grammar. Faithfulness is not a property of a language; it is a property of a
construct, a lowering, and a deployment. A capability table derived from a
specification is therefore not a capability declaration, and treating it as one
is precisely the over-declaration error of \cref{rem:honesty-assumption}.

\begin{remark}[Where \cref{asm:independence} fails]
\label{rem:independence-fails}
The assumption fails when two constructs interact. In SPARQL, a path expression
inside a pattern and an inline value block are formally independent positions,
but the \emph{expansion} of an abbreviated object list happens during
translation to the abstract algebra, before evaluation, and interacts with
path expressions in a way that surprises authors --- as \cref{sec:interp}
documents. Where independence fails, \cref{thm:lowering} does not apply and the
lowering must be verified for the interacting pair jointly. The prototype
handles this by never emitting the abbreviated form (\cref{cons:longhand}),
which restores independence by removing one of the interacting constructs from
the emitted language altogether.
\end{remark}

\section{The intermediate representation}
\label{sec:ir}

\subsection{Steps}

\begin{definition}[Step]
\label{def:step}
A \emph{step} is a tuple
$\Step = (x, \Src, \rho, \beta, b)$
where $x$ is the variable the step binds, $\Src$ is a source, $\rho$ is an
abstract request, $\beta = (y_1,\dots,y_k)$ is a tuple of variables whose
values are supplied to $\rho$ as input, and $b \in \Rset_{>0} \cup \{\infty\}$
is the step's budget annotation.
\end{definition}

The shape $\Step = \text{source} \times \text{request} \times \text{binding}$
is the whole of the intermediate representation, and it is the reason a source
adapter is small: of the five components, only $\Low_\Src$ and $\Ext_\Src$ are
source-specific, and both are functions of the abstract request alone.

\begin{definition}[Plan]
\label{def:plan}
A \emph{plan} $\Plan$ is a finite sequence $\Step_1,\dots,\Step_m$ of steps
together with a total budget $\Bud$, such that the variable bound by $\Step_i$
is distinct from that bound by $\Step_j$ for $i \neq j$, and every variable in
$\beta_i$ is bound by some $\Step_j$ with $j < i$. The \emph{dependency graph}
$G(\Plan)$ has the steps as vertices and an edge $\Step_j \to \Step_i$ whenever
$y \in \beta_i$ is bound by $\Step_j$.
\end{definition}

By construction $G(\Plan)$ is a directed acyclic graph and the sequence order
is one of its topological orders. Nothing below depends on which one.

\begin{example}[The motivating question as a plan]
\label{ex:plan}
\begin{lstlisting}[caption={The question of \cref{sec:intro} as a plan. Compare \cref{lst:script}: no query text, no batching constant, explicit budget, explicit failure handling.},label=lst:plan]
plan enzymes_in_shared_pathways {
  budget 400 requests

  let acids = from CHEBI
      ask descendants_of("CHEBI:35238")
      within 20

  let rxns  = from RHEA
      ask reactions_consuming(?c)
      with ?c in acids
      within 120
      else fail unresolved

  let ecs   = from RHEA
      ask enzyme_of(?r)
      with ?r in rxns
      within 60

  let kids  = map ecs via ec_to_kegg
      expect partial 0.6

  let paths = from KEGG
      ask link("pathway", ?k)
      with ?k in kids
      within 200
      when starved emit partial

  emit paths with provenance
}
\end{lstlisting}
The \texttt{expect partial 0.6} annotation on the translation step is not
decoration: it is the assertion checked by \cref{def:retention-check}, and a
plan whose translation retains less than the asserted fraction stops rather
than proceeding on a silently decimated input --- which is defect (D3) of
\cref{sec:intro}, caught at the step that causes it rather than at the step that
suffers it.
\end{example}

\subsection{The static capability check}

\begin{definition}[Well-capability]
\label{def:wellcap}
A step $\Step = (x,\Src,\rho,\beta,b)$ is \emph{well-capability} if
$\Req(\rho) \subseteq \Capset(\Src)$, and additionally
$\textsf{batch} \in \Capset(\Src)$ whenever $\beta \neq ()$ and the step is
annotated for batched execution. A plan is well-capability if every step is.
\end{definition}

\begin{theorem}[Static decidability of compilability]
\label{thm:static}
Let $\Plan$ be a plan over sources with honest declarations satisfying
\cref{asm:independence}. Then:
\begin{enumerate}[label=(\alph*),leftmargin=2.4em]
\item whether $\Plan$ is well-capability is decidable in time
$O\bigl(m \cdot |\Feat|\bigr)$, where $m$ is the number of steps, without
issuing any request;
\item if $\Plan$ is well-capability then every $\Low_{\Src_i}(\rho_i)$ is
defined and faithful;
\item if $\Plan$ is not well-capability then there is a step $\Step_i$ and a
capability $f \in \Req(\rho_i) \setminus \Capset(\Src_i)$ such that no execution
of $\Plan$ can produce a faithful result for $\Step_i$, and the pair
$(\Step_i, f)$ is computable within the same bound.
\end{enumerate}
\end{theorem}

\begin{proof}
(a) $\Req(\rho_i)$ is determined syntactically from $\rho_i$ by structural
recursion over the abstract request, and $\Capset(\Src_i)$ is a declared finite
set; the containment test is a subset test over a fixed finite vocabulary,
costing $O(|\Feat|)$ per step with bitset representation. There are $m$ steps
and no step's test depends on any other step's result, so the total is
$O(m|\Feat|)$. Nothing in the test consults $\eta_{\Src_i}$, so no request is
issued.

(b) Immediate from \cref{thm:lowering} applied stepwise, the hypotheses of
which hold by assumption.

(c) If $\Plan$ is not well-capability then some step fails the containment
test; take the least such $i$ in the sequence order and any witness
$f \in \Req(\rho_i)\setminus\Capset(\Src_i)$, both produced by the same scan. By
\cref{thm:lowering}, $\Low_{\Src_i}(\rho_i)$ is either undefined or unfaithful.
In the first case no request is ever issued for $\Step_i$ and it has no result.
In the second, a result is produced but differs from
$\llbracket \rho_i \rrbracket_D$ on some dataset, so it is not faithful. Since
the executor's behaviour at $\Step_i$ depends on the source only through
$\Low_{\Src_i}$ and $\Ext_{\Src_i}$, no scheduling choice, retry, or budget
allocation changes this.
\end{proof}

\begin{corollary}[Refusal before contact]
\label{cor:refuse-before-contact}
A plan that is not well-capability can be refused, with the offending
(step, capability) pair named, before any network request is issued.
\end{corollary}

\begin{proof}
Immediate from \cref{thm:static}(a) and (c): the test is decidable without
issuing a request, and on failure it computes the witness.
\end{proof}

\Cref{cor:refuse-before-contact} is \cref{prin:refusal} made operational, and
it is the concrete payoff of separating $\Capset$ from $\Ext$ in
\cref{def:source}. It is worth being precise about what it buys. It does not
guarantee that a well-capability plan succeeds --- \cref{sec:verdicts} lists
four ways it can fail with every capability present, and
\cref{prop:necessary-not-sufficient} adds a fifth. It guarantees only that one
particular family of failures, the family in which a plan asks a source for
something it cannot do, is converted from a runtime symptom into a
compile-time diagnosis naming its own cause.

\begin{proposition}[The check is not conservative in the useless direction]
\label{prop:check-tight}
There is no plan that is refused by \cref{def:wellcap} and that could
nonetheless be executed faithfully against sources with honest declarations.
\end{proposition}

\begin{proof}
Suppose $\Plan$ is refused. Then some step has
$f \in \Req(\rho_i)\setminus\Capset(\Src_i)$, and by \cref{thm:static}(c) no
execution produces a faithful result for that step. So the refusal excludes no
faithfully executable plan.
\end{proof}

\Cref{prop:check-tight} is the converse guarantee to
\cref{cor:refuse-before-contact} and is what distinguishes a capability check
from a conservative approximation: the check refuses exactly the plans that
cannot work, given honest declarations. Under \emph{dishonest} declarations it
inherits the dishonesty in both directions, which is
\cref{rem:honesty-assumption} again and is the reason that remark, rather than
this proposition, is the honest summary of what has been established.
